\documentclass[11pt,a4paper]{article}
\usepackage[T1]{fontenc}
\usepackage{lmodern,amsmath,amssymb,amsthm,mathtools,microtype}
\usepackage[margin=27mm]{geometry}
\usepackage[colorlinks=true,linkcolor=blue,citecolor=blue,urlcolor=blue]{hyperref}
\setlength{\emergencystretch}{2em}
\newtheorem{theorem}{Theorem}[section]
\newtheorem{lemma}[theorem]{Lemma}
\newtheorem{proposition}[theorem]{Proposition}
\newtheorem{definition}[theorem]{Definition}
\newtheorem{corollary}[theorem]{Corollary}
\newcommand{\Q}{\mathbb Q}\newcommand{\Z}{\mathbb Z}
\DeclareMathOperator{\Tr}{Tr}\DeclareMathOperator{\TC}{TC}
\title{One actual-trace square in a derivative detector}
\author{}\date{17 November 2026}
\begin{document}\maketitle

\begin{abstract}
The ordinary-$abc$ derivative detector admits a nine-coordinate
sound source with one finite quadratic test at its actual reciprocal
trace. Its completeness follows from an additional witness
requirement, for which the existing preparation supplies several local
conditions but no global rational realization. We verify the full
guard and terminal-root hypotheses. The source is invariant under
simultaneous negation of the input and its two associated elliptic
points. Consequently actual-trace square witnesses for all sufficiently
large positive integers would suffice for an integer definition in
nine variables, after multiplication by finitely many integer input
factors. Neither that witness assertion nor a new defining count
is proved.
\end{abstract}

\section{The retained source}

Use the nine-shift derivative source of \cite{Derivative}, including
its arithmetic hypotheses: ordinary rational $abc$, the retained
elliptic group and point charts, formal comparison, trace heights,
and prime preparation. The fixed exceptional set $S$ is unchanged
and contains all primes at most $76176$.
The curve and generator are
\[
 K=\Q(\sqrt2),\qquad E:y^2=x^3-(11-6\sqrt2)x,
\]
\[
 P=\left(\frac{57-28\sqrt2}{8},\frac{-660+417\sqrt2}{32}\right),
 \qquad 2E(K)=\Z P.
\]
For the three point roles $A,H,U$, put $f=1/x$, $\zeta=-x/y$, and
\[
 r=\frac{\zeta(U)}{\zeta(H)},\qquad
 r'=\overline{\frac{\zeta(U-A)}{\zeta(H-P)}}.
\]
The six tests are
\[
 \begin{gathered}
 C_1=f(H),\quad C_2=f(U),\quad
 C_3=\overline{f(H-P)},\quad C_4=\overline{f(U-A)},\\
 C_5=r-r',\qquad C_6=r-X.
 \end{gathered}
\]
The point chart for each role retains two rational coordinates and
its original free rank-four point algebra.

Fix an odd integer $B\ge3$ and define
\begin{equation}\label{eq:source}
\begin{gathered}
 \beta=X^2+B,\quad T=\Tr x(A),\quad M=\beta^4+T^4,\quad
 a=M^4,\quad t=M^8-z^2,\\
 F(a,z)=\prod_{c\in\{-10,-8,-5,1,2,3,4,6,7\}}
                  ((a+c)^2-z^2),\qquad u=-18\frac{F}{F_a},\\
 D_i=\frac{C_i^2}{1+uC_i^2},\qquad
 W=\prod_{i=1}^6\left(1+\frac{D_i^2}{u}\right),\qquad
 V=\Tr(W^{-1}),\\
 w=\frac{u}{(1+u)^2},\qquad b=\frac{4w}{V^2}.
\end{gathered}
\end{equation}
Only $z$ is an unrestricted slope coordinate. The integral circuit
uses the equivalent quotient $u=-9\mathcal G/\mathcal H$, where
$\mathcal G(A,Z)=\prod_c((A+c)^2-Z)$ and the primitive polynomial
$\mathcal H=\mathcal G_A/2$ is evaluated at $(a,z^2)$.

Retain every original point, chart, affine-difference and quotient
guard. Also retain the nonzero conditions on $M,t,F,F_a,u,1+u$,
the six balanced denominator norms, the two-slot product norm,
and the actual trace $V$. The existing nonzero guard on $7b$ may
be retained. Every syntactic rational denominator is guarded,
and all six preceding quadratic-stage radicands in the point algebra
remain units. No individual $C_i$ is required to be nonzero.

Use the actual-trace test of \cite{Conic} in place of the two
unrestricted intrinsic-pole parameters and their finite ordinate:
\begin{equation}\label{eq:conic}
                  Y^2=\mathcal C,\qquad
                  \mathcal C=V^2+w(2-V).
\end{equation}
The root $Y$ is a finite generator. Its value and its radicand may
vanish in soundness; no free conic parameter is introduced.

\section{The local implication and preparation}

\begin{lemma}\label{lem:normalized}
Every odd valuation of $w=u/(1+u)^2$ is positive. If
\eqref{eq:conic} has a rational solution, then $b$ has no negative
odd valuation outside $S$.
\end{lemma}
\begin{proof}
For $\nu=v_p(u)$ one has
\[
 v_p(w)=
 \begin{cases}
 \nu,&\nu>0,\\
 -\nu,&\nu<0,\\
 -2v_p(1+u),&\nu=0.
 \end{cases}
\]
Suppose $v_p(b)<0$ is odd, with $p\notin S$.
Then $v_p(w)$ is odd and positive, and
$2v_p(V)>v_p(w)$. Thus $V$ has positive valuation and $2-V$
is a unit. The two terms of $\mathcal C$ have distinct valuations,
so $v_p(\mathcal C)=v_p(w)$ is odd. In particular
$\mathcal C\ne0$ and cannot be a rational square. This includes
a zero putative terminal ordinate. Outside $S$, the conclusion is
exactly the implication formerly supplied by the intrinsic equation
for $7b$.
\end{proof}

\begin{proposition}\label{prop:local}
The following assertions hold for the actual derivative detector.
\begin{enumerate}
\item At every selected contact of the retained preparation,
$\mathcal C$ is a nonzero local square.
\item At every prime outside $S$ where $v_p(u)<0$ is odd,
$\mathcal C$ is a nonzero local square, including new primes of
$F_a$.
\item The fixed finite-place preparation can make $\mathcal C$
a nonzero local square at each prescribed place, including two.
\item At every prepared real point-algebra character,
$\mathcal C>0$ and all new denominator and trace guards are nonzero.
\end{enumerate}
\end{proposition}
\begin{proof}
At a selected contact, $v_p(u)=v_p(w)=1$ and $V$ is a unit,
congruent to one or two. Hence
\begin{equation}\label{eq:local-square}
          \frac{\mathcal C}{V^2}
               =1+\frac{w(2-V)}{V^2}\in1+p\Z_p,
\end{equation}
and Hensel's lemma applies.

For any odd $\nu=v_p(u)$, the balanced identity gives
\[
                  v(D_i)=\max(2v(C_i),-\nu)
\]
at either place of the unramified quadratic algebra. Zero tests have
valuation infinity. If $\nu<0$, then
$v(D_i^2/u)\ge-3\nu>0$ for every test. Both slots of $W$ are
one modulo $p$, so $V\equiv2\pmod p$, while
$v_p(w)=-\nu>0$. Equation~\eqref{eq:local-square} proves the
second assertion.

At each fixed preparation place the source asymptotics are
\[
                      u\sim z^2/a,\qquad w\sim a/z^2,
                         \qquad V\longrightarrow2.
\]
The already chosen nonzero test valuations stay fixed; identically
zero tests stay zero. Taking the existing preparation sufficiently
deep puts $w(2-V)/V^2$ in $p\Z_p$ at an odd place and in
$8\Z_2$ at two. This proves the third assertion using the same point family and the single source
slope, without an independent trace parameter or a change in the
order of preparation.

At every prepared real character $a>10$ and $z>a+10$. Thus
$u,w>0$, every denominator $1+uC_i^2$ is positive in both slots,
and every slot of $W$ is at least one. Consequently
\begin{equation}\label{eq:positive}
                   0<V\le2,\qquad
                    \mathcal C=V^2+w(2-V)>0.
\end{equation}
All corresponding denominator, product and trace norms are nonzero.
The original point and difference guard preparation is retained.
\end{proof}

These local statements do not settle squareness at every new prime
of even physical valuation. They do not imply rational squareness.

\section{The exact witness requirement}

\begin{definition}\label{def:TC}
For a fixed odd $B\ge3$ and a nonzero integer $n$, write
$\TC_B(n)$ if there exist the three actual rational point-chart
pairs and one rational slope $z$, with $X=n$ and $A=[n]P$,
satisfying every retained guard and the whole-character preparation
of \cite{Derivative}, for which the actual value
\[
                          V^2+\frac{u}{(1+u)^2}(2-V)
\]
is a rational square. The preparation includes the original point
and difference avoidance conditions, fixed-place and selected-contact
conditions, and $a>10$, $z>a+10$ at every real character.
The multiplication family $U=[n]H$ is permitted.

Write $\TC(B)$ if $\TC_B(n)$ holds for every nonzero integer $n$.
Write $\TC^+_{\mathrm{cof}}(B)$ if some integer $N\ge1$ has
$\TC_B(n)$ for every integer $n\ge N$.
\end{definition}

These conditions concern the existing point tuple and the same slope
appearing in $F$. They do not replace $V$ by a freely chosen rational
trace, or introduce an extra norm, conic, multiplier or square-class
coordinate. The value $B$ must exceed the same ineffective soundness
threshold as in the derivative construction.

\begin{theorem}\label{thm:transfer}
Under the retained arithmetic premises and ordinary rational $abc$,
there is an ineffective $B_0$ and a computable family $P_B$ of
integer polynomials in nine total variables, including the input,
with these properties.
\begin{enumerate}
\item For every odd $B\ge B_0$, the input projection of $P_B=0$
is contained in $\Z$ and contains zero.
\item For every such $B$, each input satisfying $\TC_B(n)$ is
in that projection. In particular $\TC(B)$ implies that the
projection is exactly $\Z$.
\item The native finite algebra has rank $128=4^3\cdot2$.
\end{enumerate}
No assertion that any $B\ge B_0$ satisfies $\TC(B)$ is made.
\end{theorem}
\begin{proof}
Let $v_p(u)=\nu>0$ be odd. Every factor of $W$ has valuation
$\min(0,2v(D_i)-\nu)\le0$. Let $q\le0$ be the largest
valuation of $W$ over the places above $p$. Then $v_p(V)\ge-q$.
Lemma~\ref{lem:normalized} implies
\[
                         1\le v_p(b)\le\nu+2q.
\]
At a place attaining $q$, every individual factor has valuation
at least $q> -\nu/2$, so $v(D_i)>\nu/4>0$. The balanced
identity gives $v(C_i)>0$ for all six original tests.
The retained formal comparison therefore implies $p$-integrality
of $X$ and $X\equiv n\pmod p$, where $A=[n]P$.
Thus the derivative source's extraction and height argument apply
unchanged. Its strict absorption margin is
\[
 \frac{297}{200}
 -\frac{501}{997}\left(\frac{201}{100}+\frac{1001}{1994}\right)
                 =\frac{44272779}{198801800}>0.
\]
It follows that every guarded rational source solution has $X=n$
for sufficiently large fixed odd $B$. No new exceptional primes or
height-support losses are introduced.

We verify the finite algebra and every new denominator.
Write $W=A_0+B_0\sqrt2$ and $N_W=A_0^2-2B_0^2$. Then
\begin{equation}\label{eq:cleared}
 \mathcal C=
           \frac{4A_0^2+2wN_W(N_W-A_0)}{N_W^2}.
\end{equation}
There is no denominator beyond those of the guarded weight and
detector norm, and no division by $2-V$. The original guard on
$V$ remains necessary to define $b$.

If $\mathcal C=N_C/D_C$, adjoin a monic root $q^2=N_CD_C$.
If a syntactic square is extracted, write $N_CD_C=S_0^2c$ and
adjoin $q_0^2=c$; then $Y=S_0q_0/D_C$. Soundness requires
no inverse of $S_0$. The original free rank-64 point algebra and
this one quadratic step have rank 128. Every preceding radicand
remains a unit; only the last one may vanish. The subalgebra norms
of \cite{Subalgebra} and terminal-root separator of \cite{Terminal}
therefore apply with their stated hypotheses. They use one rational
coordinate and force a rational geometric character of the complete
source, even on a nonreduced terminal fiber. The pole implication
of Lemma~\ref{lem:normalized} remains valid there.

If $\TC_B(n)$ holds, its selected rational character extends to
the new quadratic algebra because the actual $\mathcal C$ is square.
All whole-character guards hold, and \eqref{eq:positive} makes
every radicand value nonzero on the prepared locus. Consequently
any extracted $S_0$ is automatically nonzero there, so the chosen
rational ordinate gives the reduced quadratic root as well. No
additional witness-avoidance parameter is needed. The separator is
nonzero at every geometric character and supplies its single rational
inverse coordinate. Nonreal terminal roots away from the prepared
locus do not affect descent.

The unrestricted coordinates are
\[
 1\text{ input}+6\text{ point coordinates}+1\text{ slope}
                         +1\text{ separator}=9.
\]
Clearing the guarded denominators and taking the specified norm
determinant gives a computable integer family. Multiplication by $X$
includes zero. No degree or coefficient bound is asserted here.
\end{proof}

\section{Sign symmetry and finite exceptions}

\begin{proposition}\label{prop:sign}
For every fixed $B$ and nonzero integer $n$,
\[
                              \TC_B(n)\iff\TC_B(-n).
\]
\end{proposition}
\begin{proof}
Apply
\begin{equation}\label{eq:sign}
                     (X,A,H,U,z)\longmapsto(-X,-A,H,-U,z).
\end{equation}
Point abscissas and $T$ remain fixed, so all scalar functions
$\beta,M,a,t,F,F_a,u,w$ are fixed. Since $f(-Q)=f(Q)$ and
$\zeta(-Q)=-\zeta(Q)$, the first four tests remain fixed while
$r,r',C_5,C_6$ negate. Their squares, and hence $D_i,W,V,b$ and
$\mathcal C$, remain fixed.

The retained guards are preserved as well. In the source's homogeneous
difference formula, with $\Delta x=x_1-x_2$ and $D=11-6\sqrt2$,
\[
 \begin{aligned}
 X_{12}&=(x_1+x_2)(x_1x_2-D)+2y_1y_2,\\
 Y_{12}&=(y_1x_2+y_2x_1)(\Delta x)^2-(y_1+y_2)X_{12}.
 \end{aligned}
\]
Simultaneously negating both points fixes $\Delta x,X_{12}$ and
negates $Y_{12}$. Thus the difference-denominator norms and all
quotient guards stay nonzero. The difference $H-P$ is unchanged.
The original point inverses have unchanged norms. No comparison
test is separately guarded, and all added detector guards depend
on the unchanged squares or scalar functions.

In each of the $A,U$ rank-four blocks, the finite generators
$(\mathsf U,\mathsf V,\mathsf W)$ transform as
$(-\mathsf U,-\mathsf V,\mathsf W)$. This preserves the point
multiplication relations and permutes all geometric characters.
Therefore the preceding radicand-unit conditions and every
whole-character guard are preserved. The real inequalities and
all local valuation and congruence preparations are unchanged after
the corresponding point signs are reversed. The actual square in
\eqref{eq:conic} is unchanged. Applying the involution twice proves
the equivalence.
\end{proof}

\begin{corollary}\label{cor:cofinite}
Suppose the retained arithmetic premises and ordinary rational $abc$
hold, and some odd $B\ge B_0$ satisfies
$\TC^+_{\mathrm{cof}}(B)$. Then the integers have a Diophantine
definition with nine total variables. More precisely, if
$\TC_B(n)$ holds for all $n\ge N\ge1$, the polynomial
\[
              \widetilde P_B
                =P_B\prod_{m=1}^{N-1}(X^2-m^2)
\]
has input projection exactly $\Z$.
\end{corollary}
\begin{proof}
Proposition~\ref{prop:sign} and Theorem~\ref{thm:transfer} supply
every integer with $|n|\ge N$. The factor $X$ already supplies
zero, and the displayed factors supply the remaining integers.
Every new root has integer input, so soundness is preserved.
No variable or inverse is added.
\end{proof}

The hypothesis of Corollary~\ref{cor:cofinite} is unproved. The
integer $N$ and sufficient shift $B$ are not numerically specified.
The result supplies no effective exception set or bound on the
extra degree and coefficients. Its significance is that an eventual
positive-input realization theorem suffices; completeness at finitely
many exceptional inputs can be included without increasing the
number of variables.

\section{Exact algebraic verification}

\begingroup\raggedright
The checker
\nolinkurl{computation/pole_parameters/single_trace_v64/check.py}
verifies the cleared trace identity, invariance of the homogeneous
difference formulas, and the sign action on all six rank-four point
relations. These are exact symbolic controls of the added identities.
The valuation, height, finite-algebra and cofinite implications are
proved above; no global actual-trace realization is inferred from
those calculations.\par\endgroup

\begin{thebibliography}{9}\raggedright
\bibitem{Derivative}\emph{An eleven-variable construction under rational
$abc$}, repository manuscript, 17 November 2026, version 63,
\nolinkurl{papers/conditional_few_variables/derivative_detector11_20261117_v63.tex},
published source commit
\nolinkurl{f70057ffce6b8318a79a7e3648779067fd98c2e3}.
\bibitem{Conic}\emph{A rational conic criterion using the actual detector
trace}, repository manuscript, 17 November 2026, version 57,
\nolinkurl{papers/pole_parameters/trace_conic_criterion_20261117_v57.tex}.
\bibitem{Subalgebra}\emph{Smaller defining polynomials from subalgebra
norms}, repository manuscript, version 30,
\nolinkurl{papers/finite_algebra/subalgebra_norms_20261114_v30.tex}.
\bibitem{Terminal}\emph{Smaller norm polynomials with a possibly singular
final quadratic extension}, repository manuscript, version 40,
Section 3,
\nolinkurl{papers/finite_algebra/terminal_quadratic_root_20261115_v40.tex}.
\end{thebibliography}
\end{document}
